% !TeX program = xelatex
\documentclass[10pt,a4paper]{ctexart} % 设置文章类型，正文字号与纸张大小

% 页面与正文常用宏包
\usepackage[margin=2cm]{geometry} % 设置页面四周页边距
\usepackage{graphicx}             % 插入图片
\usepackage{amsmath}              % 提供cases、aligned等数学公式环境
\usepackage{amssymb}              % 提供\therefore等扩展数学符号
\usepackage{booktabs}             % 提供\toprule、\midrule、\bottomrule三线表命令
\usepackage{float}                % 提供[H]固定图片、表格位置
\usepackage{subcaption}           % 提供subfigure子图环境及子图标题
\usepackage{wrapfig}              % 提供wrapfigure图文环绕环境
\usepackage{fancyhdr}             % 自定义页眉和页脚
\usepackage{hyperref}             % 生成超链接

% 页眉页脚
\pagestyle{fancy}                  % 启用fancyhdr定义的页面样式
\fancyhf{}                         % 清空默认页眉和页脚
\fancyfoot[C]{\thepage}            % 在页脚中央显示页码
\renewcommand{\headrulewidth}{0pt} % 隐藏页眉横线
\renewcommand{\footrulewidth}{0pt} % 隐藏页脚横线

% 超链接
\hypersetup{
  colorlinks=true,                % 使用彩色文字链接，而非带边框的链接
  linkcolor=red,                 % 目录、公式及图表交叉引用的颜色
  filecolor=cyan,                 % 本地文件链接的颜色
  urlcolor=blue,                  % 网址链接的颜色
  pdftitle={物理实验预习报告模版},     % 写入PDF文档属性的标题
  bookmarksnumbered=true          % PDF书签显示章节编号
}

% 定义字号命令
\newcommand{\chuhao}{\fontsize{42pt}{63pt}\selectfont}
\newcommand{\xiaochu}{\fontsize{36pt}{54pt}\selectfont}
\newcommand{\yihao}{\fontsize{26pt}{39pt}\selectfont}
\newcommand{\xiaoyi}{\fontsize{24pt}{36pt}\selectfont}
\newcommand{\erhao}{\fontsize{22pt}{33pt}\selectfont}
\newcommand{\xiaoer}{\fontsize{18pt}{27pt}\selectfont}
\newcommand{\sanhao}{\fontsize{16pt}{24pt}\selectfont}
\newcommand{\xiaosan}{\fontsize{15pt}{22.5pt}\selectfont}
\newcommand{\sihao}{\fontsize{14pt}{21pt}\selectfont}
\newcommand{\xiaosi}{\fontsize{12pt}{18pt}\selectfont}
\newcommand{\wuhao}{\fontsize{10.5pt}{15.75pt}\selectfont}
\newcommand{\xiaowu}{\fontsize{9pt}{13.5pt}\selectfont}

% 章节格式
\ctexset{
  section={
    format={\CJKfontspec[AutoFakeBold=2.5]{KaiTi}\sanhao\bfseries}, % 一级标题使用三号加粗楷体
    name={,、},                                                    % 在中文章节序号后添加顿号
    number=\chinese{section},                                      % 一级标题编号使用“一、二、三”形式
    aftername={}                                                   % 取消编号与标题文字之间的空格
  },
  subsection={
    format={\CJKfontspec[AutoFakeBold=2.5]{KaiTi}\sihao\bfseries}, % 二级标题使用四号加粗楷体
  },
  subsubsection={
    format={\CJKfontspec[AutoFakeBold=2.5]{KaiTi}\xiaosi\bfseries}, % 三级标题使用小四加粗楷体
  }
}

\begin{document}

% 封面
\thispagestyle{fancy} % 封面保留底部居中的页码
\fancyfoot[C]{\thepage}
\begin{center}
  \CJKfontspec[AutoFakeBold=2.5]{KaiTi}\bfseries % 封面中文使用楷体并人工加粗，西文同时加粗
  \vspace*{2.0cm}
  \xiaochu\textbf{物理实验预习报告}

  \vspace{3.8cm}
  \begin{minipage}{10.0cm}
    \sanhao
    \makebox[2.6cm][l]{\textbf{实验名称：}}
    \underline{\makebox[7.0cm][c]{}}\par\vspace{12pt}

    \makebox[2.6cm][l]{\textbf{指导老师：}}
    \underline{\makebox[7.0cm][c]{}}\par\vspace{12pt}
  \end{minipage}

  \vspace{3.8cm}
  \begin{minipage}{10.0cm}
    \sihao
    \makebox[1.4cm][l]{\textbf{班级：}}
    \underline{\makebox[8.0cm][c]{}}\par\vspace{10pt}

    \makebox[1.4cm][l]{\textbf{姓名：}}
    \underline{\makebox[8.0cm][c]{IllusoryRain}}\par\vspace{10pt}

    \makebox[1.4cm][l]{\textbf{学号：}}
    \underline{\makebox[8.0cm][c]{}}
  \end{minipage}

  \vspace{1.8cm}
  \sihao
    \textbf{实验日期：}
    \underline{\makebox[0.7cm][c]{}}年
    \underline{\makebox[0.7cm][c]{}}月
    \underline{\makebox[0.7cm][c]{}}日
    星期\underline{\makebox[0.7cm][c]{}}上午

  \vspace{0.7cm}
  \sihao 
    浙江大学物理实验教学中心

  \vspace{0.2cm}
\end{center}

\newpage
\xiaosi\songti\mdseries % 后续正文统一使用小四不加粗的宋体

\section{预习报告（10分）}
（注：将已经写好的“物理实验预习报告”内容拷贝过来）
  \subsection{实验综述（5分）}
  （自述实验现象、实验原理和实验方法，包括必要的光路图、电路图、公式等。不超过500字。）
  \subsection{实验重点（3分）}
  （简述本实验的学习重点，不超过100字。）
  \subsection{实验难点（2分）}
  （简述本实验的实现难点，不超过100字。）

\vspace*{\fill} % 使下面内容在页面底部显示
\noindent
\begin{minipage}{\textwidth}
  注意事项：

  \begingroup
    \setlength{\leftskip}{2em} % 1--4整体向右缩进两个汉字宽度
    1. 用PDF格式上传“实验报告”，文件名：学生姓名+学号+实验名称+周次。

    2. “实验报告”必须递交在“学在浙大”本课程内对应实验项目的“作业”模块内。

    3. “实验报告”成绩必须在“浙江大学物理实验教学中心网站”-“选课系统”内查询。

    4. 教学评价必须在“浙江大学物理实验教学中心网站”-“选课系统”内进行，学生必须进行教学评价，才能看到实验报告成绩，教学评价须在本次实验结束之后3天内进行。\par
    \endgroup

  \begin{center}
    {\CJKfontspec[AutoFakeBold=2.5]{SimSun}\bfseries 浙江大学物理实验教学中心制}
  \end{center}
\end{minipage}

\end{document}
